Un dels mètodes emprats per a detectar exoplanetes (planetes extrasolars) és l'observació del \textit{trànsit planetari}, un fenomen astronòmic que s'esdevé quan un planeta passa per davant de l'estel al voltant del qual orbita i que es percep des de la Terra per la disminució de la llum de l'estel.

El gràfic següent mostra la variació de lluminositat provocada pel trànsit d'un planeta que descriu una òrbita circular al voltant d'un estel. Aquest estel té una massa pràcticament idèntica a la massa del Sol. Considereu que la constant de Kepler d'aquest sistema és igual a la del Sistema Solar.

\figura[0.75]{fig1.png}

\begin{apartat}{1,25}
Calculeu el període i el radi de l'òrbita.
\end{apartat}

\begin{apartat}{1,25}
Determineu el mòdul de la velocitat i l'acceleració centrípeta del planeta.
\end{apartat}

\dades{Radi orbital mitjà de la Terra $= 1,00\ \mathrm{ua} = 1,50\times 10^{11}\ \mathrm{m}$.}
